\section*{Item 1}\label{it:1}

\textbf{Enunciado}: Um conversor Flyback possui os seguintes parâmetros: $V_i = 36$ V, $D = 0,4$, $N_p/N_s = 2$, $R = 20~\Omega$, $L_m = 100~\mu\text{H}$ e $C = 50~\mu\text{F}$. A frequência de comutação é de $100$ kHz. Determine:

\medskip
\textbf{a)} a tensão de saída;

\medskip
\textbf{Resolução}: inicialmente, verifica-se o modo de operação do conversor. Para isso, calcula-se a indutância de magnetização mínima necessária para operação em CCM~\cite{flyback_notas}
\begin{equation*}
	L_{m,\min}=\frac{R}{2N^2}(1-D)^2T_s,
\end{equation*}
em que $N=N_s/N_p=0{,}5$. Logo,
\begin{equation*}
	\begin{aligned}
		L_{m,\min}
		&=\frac{20}{2(0{,}5)^2}(0{,}6)^21\cdot10^{-5}\\
		&=225~\mu\mathrm{H}.
	\end{aligned}
\end{equation*}

Como o circuito possui $L_m=100~\mu$H, o conversor opera em DCM. Nesta condição, a característica de conversão é dada por
\begin{equation*}
	\frac{V_o}{V_i}
	=
	D\sqrt{\frac{RT_s}{2L_m}},
\end{equation*}
resultando em
\begin{equation*}
	\begin{aligned}
		V_o
		&=36\cdot0{,}4
		\sqrt{\frac{20(1\cdot10^{-5})}{2(1\cdot10^{-4})}}\\
		&\approx14{,}4~\mathrm{V}.
	\end{aligned}
\end{equation*}

\medskip
\textbf{b)} as correntes média, máxima e mínima no indutor;

\medskip
\textbf{Resolução}: a corrente média na indutância de magnetização é obtida por~\cite{flyback_notas}
\begin{equation*}
	\begin{aligned}
		I_{Lm,\mathrm{med}}
		&=
		\frac{V_o^2}{V_iDR},\\
		&=
		\frac{14{,}4^2}{36\cdot0{,}4\cdot20}
		=0{,}72~\mathrm{A}.
	\end{aligned}
\end{equation*}

O ripple de corrente na indutância é~\cite{flyback_notas}
\begin{equation*}
	\Delta i_{Lm}
	=
	\frac{V_iDT_s}{L_m}
	=
	\frac{36\cdot0{,}4\cdot1\cdot10^{-5}}
	{100~\mu\mathrm{H}}
	=
	1{,}44~\mathrm{A}.
\end{equation*}

Assim,~\cite{flyback_notas}
\begin{equation*}
	I_{Lm,\max}
	=
	I_{Lm,\mathrm{med}}
	+
	\frac{\Delta i_{Lm}}{2}
	=
	1{,}44~\mathrm{A},
\end{equation*}
e
\begin{equation*}
	I_{Lm,\min}
	=
	I_{Lm,\mathrm{med}}
	-
	\frac{\Delta i_{Lm}}{2}
	=
	0~\mathrm{A},
\end{equation*}
confirmando a operação em modo descontínuo.

\medskip
\textbf{c)} a ondulação (ripple) da tensão de saída;

\medskip
\textbf{Resolução}: a ondulação da tensão de saída pode ser estimada por~\cite{flyback_notas}
\begin{equation*}
	\Delta V_o
	=
	\frac{V_o}{RC}DT_s.
\end{equation*}

Substituindo os parâmetros do circuito,
\begin{equation*}
	\begin{aligned}
		\Delta V_o
		&=
		\frac{14{,}4}
		{20\cdot5\cdot10^{-5}}
		0,4\cdot1\cdot10^{-5}\\
		&\approx57{,}6~\mathrm{mV}.
	\end{aligned}
\end{equation*}

\medskip
\textbf{d)} o valor da carga $R$ no limiar entre o MCC e o MCD.

\medskip
\textbf{Resolução}: isolando $R$ da condição limite entre CCM e DCM~\cite{flyback_notas}, obtém-se
\begin{equation}
	R
	=
	\frac{2N^2L_m}
	{(1-D)^2T_s}.
\end{equation}

Substituindo os valores do circuito,
\begin{equation}
	\begin{aligned}
		R
		&=
		\frac{2(0,5)^2(1\cdot10^{-4})}
		{(0,6)^2(1\cdot10^{-5})}\\
		&\approx13{,}89~\Omega.
	\end{aligned}
\end{equation}

Como o circuito utiliza $R=20~\Omega$, confirma-se a operação em DCM.






